%==============================================================================
% FICHAMENTO CRÍTICO — Deep Learning for Automated Detection of
% Periapical Lesions
% Conforme NBR 6023 (referências) e NBR 10520 (citações)
% Capítulo 9 — Lesões Periapicais, Diagnóstico por Imagem
%==============================================================================
\section{Fichamento: \citeonline{viet2024deep}}

% --- 1. IDENTIFICAÇÃO ---
\subsection{Identificação}
\begin{tabular}{ll}
\textbf{Título:} & Deep learning for automated detection of periapical
lesions using periapical radiographs: a comparative study of Faster
R-CNN and YOLOv4 \\
\textbf{Autor(es):} & Viet, Do Hoang; Son, Le Hoang; Tuyen, Do Ngoc;
Tuan, Tran Manh; Thang, Nguyen Phu; Ngoc, Vo Truong Nhu \\
\textbf{Periódico:} & Oral Surgery, Oral Medicine, Oral Pathology and Oral Radiology \\
\textbf{Ano:} & 2024 \\
\textbf{Volume:} & 138 \\
\textbf{Número:} & 5 \\
\textbf{Páginas:} & 481--498 \\
\textbf{DOI:} & \url{https://doi.org/10.1016/j.oooo.2024.01.020} \\
\textbf{Qualis:} & A2 (Odontologia) \\
\textbf{Tipo:} & Artigo original retrospectivo \\
\end{tabular}

% --- 2. OBJETIVO ---
\subsection{Objetivo}
Desenvolver e comparar duas arquiteturas de deep learning baseadas em
detecção de objetos — Faster R-CNN e YOLOv4 — para detecção e
classificação automatizada de lesões periapicais em radiografias
periapicais, utilizando o sistema de escore do Periapical Index (PAI)
para classificar a gravidade das lesões (PAI 3, 4 e 5). O estudo
visa estabelecer um sistema de suporte ao diagnóstico que não apenas
detecte a presença de periodontite apical, mas também estratifique
sua gravidade de acordo com a classificação radiográfica validada
internacionalmente, contribuindo para a tomada de decisão clínica
em endodontia.

% --- 3. METODOLOGIA ---
\subsection{Metodologia}
\textbf{Desenho do estudo:} Estudo retrospectivo de desenvolvimento e
comparação de dois modelos de deep learning para detecção de objetos
em radiografias periapicais digitais. \\
\textbf{Amostra:} 2.658 radiografias periapicais digitais (1.518 de
dentes superiores, 1.140 de inferiores) coletadas do arquivo
radiográfico de 2 clínicas universitárias (Vietnã), abrangendo as
três regiões do arco dental: anteriores (incisivos e caninos — 892),
pré-molares (624) e molares (1.142); 60\% das radiografias
apresentavam pelo menos uma lesão periapical confirmada por
achados clínicos e radiográficos. \\
\textbf{Métricas principais:} Sensibilidade, especificidade, acurácia,
precisão, F1-score, precisão por classe PAI (3, 4, 5), mAP@0.5,
AUC-ROC, sensibilidade por região do arco dental, tempo de inferência,
coeficiente kappa de Cohen. \\
\textbf{Validação:} Divisão treino (80\%) / validação (10\%) / teste
(10\%) estratificada por tipo de dente; validação cruzada 5-fold;
comparação com diagnóstico de 2 endodontistas experientes (padrão
de referência); análise de desempenho estratificada por região do
arco e por escore PAI; teste de robustez com diferentes condições
de angulação horizontal (bissetriz vs. paralelismo).

% --- 4. RESULTADOS PRINCIPAIS ---
\subsection{Resultados Principais}
\begin{itemize}
  \item O modelo Faster R-CNN (ResNet-101 backbone) alcançou
        sensibilidade geral de 0,932 (IC 95\%: 0,914--0,948),
        especificidade de 0,917 (IC 95\%: 0,893--0,938), acurácia de
        0,926 e F1-score de 0,928 para detecção binária de
        periodontite apical (presença vs. ausência).
  \item O modelo YOLOv4 (CSPDarknet53 backbone) apresentou
        sensibilidade de 0,889, especificidade de 0,871, acurácia de
        0,882 e F1-score de 0,890, sendo significativamente inferior
        ao Faster R-CNN para todas as métricas (p $<$ 0,01, teste de
        McNemar), com diferença de 4,3 pontos percentuais na acurácia.
  \item Na classificação por escore PAI, o Faster R-CNN apresentou
        acurácia de 89,0\% para PAI 3 (lesão leve), 83,0\% para
        PAI 4 (moderada) e 91,8\% para PAI 5 (grave), enquanto o
        YOLOv4 obteve 68,1\%, 50,9\% e 65,3\% respectivamente,
        evidenciando a superioridade do Faster R-CNN especialmente
        para lesões moderadas (PAI 4), onde a classificação é
        clinicamente mais desafiadora.
  \item A sensibilidade estratificada por região do arco dental para
        o Faster R-CNN foi de 0,942 (anteriores), 0,924 (pré-molares)
        e 0,931 (molares), sem diferença estatisticamente significativa
        entre regiões (p = 0,37), indicando desempenho homogêneo
        independentemente da complexidade anatômica regional.
  \item A concordância entre o Faster R-CNN e o consenso de
        endodontistas apresentou kappa de Cohen de 0,847
        (concordância quase perfeita), superior à concordância
        interobservador entre os endodontistas (kappa = 0,812),
        sugerindo que o modelo pode atuar como padronizador
        diagnóstico.
  \item O tempo de inferência por radiografia foi de 47 ms (Faster
        R-CNN) e 23 ms (YOLOv4) em GPU NVIDIA RTX 3060, ambos
        clinicamente viáveis para uso em tempo real, com YOLOv4
        sendo 2,04$\times$ mais rápido, porém com desempenho inferior.
\end{itemize}

% --- 5. LIMITAÇÕES ---
\subsection{Limitações}
\begin{itemize}
  \item Dataset unicêntrico vietnamita: as 2.658 radiografias foram
        provenientes de duas clínicas no mesmo país, com predomínio
        de população do sudeste asiático; a morfologia dentária,
        a prevalência de lesões e os padrões de tratamento podem
        diferir significativamente de outras populações, limitando
        a generalização global.
  \item A classificação PAI foi baseada exclusivamente em radiografias
        periapicais, sem confirmação histopatológica ou microbiológica;
        o PAI é um índice radiográfico que não se correlaciona
        perfeitamente com o estado histológico da lesão, podendo
        haver discordância entre o escore radiográfico e a gravidade
        real da inflamação periapical.
  \item O desbalanceamento entre os escores PAI — PAI 3 (38,5\%),
        PAI 4 (33,1\%) e PAI 5 (28,4\%) da amostra com lesões —
        não foi corrigido com técnicas de balanceamento (SMOTE,
        weighted loss), o que pode ter impactado negativamente o
        desempenho do YOLOv4 na classe PAI 4 (mais desafiadora e com
        sobreposição de características com PAI 3 e 5).
  \item Apenas radiografias periapicais foram utilizadas; a inclusão
        de radiografias panorâmicas e CBCT poderia fornecer uma
        avaliação mais abrangente e permitir a detecção de lesões
        em dentes com sobreposição anatômica ou em estágios iniciais
        (PAI 1 e 2), que não foram incluídos no estudo.
  \item O estudo não avaliou o efeito da incorporação do modelo no
        fluxo de trabalho clínico real, como redução do tempo de
        diagnóstico, taxa de detecção de lesões incidentais ou
        impacto na decisão terapêutica (tratamento conservador vs.
        cirúrgico), limitando a evidência de utilidade clínica.
\end{itemize}

% --- 6. CONTRIBUIÇÃO PARA O LIVRO ---
\subsection{Contribuição para o Livro}
Este artigo fundamenta a Seção 9.4 do Capítulo 9 (Lesões Periapicais),
onde é utilizado como estudo de caso principal para modelos de
detecção de lesões periapicais em radiografias periapicais. A
comparação sistemática entre Faster R-CNN e YOLOv4 é explorada na
Seção 9.4.2 (Arquiteturas de detecção) para ilustrar o trade-off
entre acurácia e velocidade computacional no contexto clínico. A
classificação por escores PAI (3, 4, 5) é adotada no livro como
sistema de referência para a Seção 9.3 (Estratificação de gravidade),
onde se argumenta que a IA pode não apenas detectar lesões, mas
também auxiliar na padronização da classificação radiográfica. A
concordância quase perfeita (kappa = 0,847) e a superioridade sobre
a concordância interobservador humana são citadas na Seção 9.5 como
evidência central para a tese do livro de que sistemas de IA bem
calibrados podem reduzir a variabilidade diagnóstica em endodontia.
A limitação de ausência de lesões PAI 1 e 2 é discutida na Seção 9.7
como lacuna de pesquisa, motivando a especificação de novos datasets
com espectro completo de lesões.

% --- 7. ANÁLISE CRÍTICA ---
\subsection{Análise Crítica}
\begin{itemize}
  \item \textbf{Validade interna:} Alta -- Amostra substancial (2.658
        radiografias), comparação direta entre duas arquiteturas
        state-of-the-art, padrão de referência por consenso de
        especialistas, análise estatística robusta (IC 95\%, teste
        de McNemar, kappa, estratificação por região e PAI).
  \item \textbf{Validade externa:} Baixa a Média -- Dataset
        unicêntrico (Vietnã), sem validação externa em população
        independente; a homogeneidade populacional e a origem
        única dos exames (região geográfica restrita) limitam
        severamente a generalização para outros contextos clínicos
        e demográficos.
  \item \textbf{Reprodutibilidade:} Média -- Metodologia descrita em
        detalhe suficiente para replicação, mas código e pesos dos
        modelos não disponibilizados; as anotações PAI não estão
        publicamente acessíveis, embora os autores declarem
        disponibilidade mediante solicitação razoável.
  \item \textbf{Relevância clínica:} Alta -- A periodontite apical é
        uma das doenças mais prevalentes em endodontia (afetando
        52\% dos indivíduos em nível global segundo revisões
        sistemáticas); a detecção precoce e a classificação precisa
        da gravidade são fundamentais para o plano de tratamento e
        o prognóstico.
  \item \textbf{Conflito de interesses:} Não -- Declaração de ausência
        de conflitos; financiamento por fundações de pesquisa
        vietnamitas (NAFOSTED) e bolsas CAPES (cooperação Brasil).
  \item \textbf{Viés identificado:} Viés de seleção — a amostra
        incluiu apenas radiografias de pacientes que já haviam sido
        encaminhados para avaliação endodôntica em clínicas
        universitárias, resultando em prevalência de lesões maior
        que a população geral (60\% vs. 52\% estimado), o que pode
        superestimar o valor preditivo positivo do modelo.
\end{itemize}

% --- 8. CITAÇÕES RELEVANTES ---
\subsection{Citações Relevantes}
\begin{citacao}
``Faster R-CNN achieved a sensitivity of 0.932 and specificity of
0.917 for periapical lesion detection, significantly outperforming
YOLOv4 (0.889 and 0.871, p $<$ 0.01). For PAI scoring classification,
Faster R-CNN correctly classified 89.0\% of PAI 3, 83.0\% of PAI 4,
and 91.8\% of PAI 5 cases, demonstrating that two-stage detectors
offer superior diagnostic granularity for lesion severity
stratification in endodontics.'' (p. 491).
\end{citacao}

% --- 9. MAPEAMENTO SDD ---
\subsection{Mapeamento SDD}
\textbf{Problema:} Detecção e classificação de gravidade de lesões
periapicais em radiografias periapicais com variabilidade
interobservador significativa, limitando a padronização do
diagnóstico e o planejamento terapêutico em endodontia. \\
\textbf{Entrada:} Radiografias periapicais digitais de dentes
anteriores, pré-molares e molares (2.658 exames); 1.518 superiores,
1.140 inferiores. \\
\textbf{Saída:} Detecção binária de periodontite apical (presença/
ausência) com caixa delimitadora; classificação PAI (3, 4, 5) com
probabilidades associadas; sensibilidade de 0,932, especificidade
de 0,917, kappa de 0,847. \\
\textbf{Critérios:} Sensibilidade $>$ 0,90; especificidade $>$ 0,90;
F1-score $>$ 0,92; kappa $>$ 0,80; acurácia PAI $>$ 80\% por classe;
validação cruzada 5-fold; tempo de inferência $<$ 100 ms/imagem. \\
\textbf{Confiança:} 0,72 — Estudo com amostra robusta e comparação
metodologicamente sólida, mas severamente limitado pela ausência de
validação externa e pela natureza unicêntrica vietnamita da amostra.
A superioridade consistente do Faster R-CNN sobre YOLOv4 é bem
suportada, mas a generalização para populações não asiáticas e para
equipamentos radiográficos diversos não foi estabelecida. A
confiança é moderada e recomenda-se validação externa antes da
adoção clínica.

% --- 10. REFERÊNCIA ABNT ---
\subsection{Referência ABNT}
\begin{verbatim}
VIET, Do Hoang et al. Deep learning for automated detection of
periapical lesions using periapical radiographs: a comparative study
of Faster R-CNN and YOLOv4. Oral Surgery, Oral Medicine, Oral
Pathology and Oral Radiology, New York, v. 138, n. 5, p. 481-498,
2024. DOI: 10.1016/j.oooo.2024.01.020.
\end{verbatim}
\endinput
